\section{Conclusion}
\label{sec:ch3-conclusion}
This work extends the HCI---HBI agenda by offering both conceptual and methodological tools for examining lived experience in smart environments. This paper demonstrated how an AfI lens can enrich the design and evaluation of smart buildings by attending to the lived, affective, and interpretive dimensions of experience. Across two exploratory studies, we surfaced five \emph{Ways} in which occupants actively navigate, assess, and adapt to their environments---through bodily cues, social attunement, and ongoing sense-making. We offered six \textit{Affective Practices} and corresponding \textit{Interaction Qualities} as a vocabulary for designing interactive systems that respond to the constantly evolving nuances of everyday experience. Future work should explore how the \textit{Affective Practices} and \textit{Interaction Qualities} introduced here can be operationalised in real-world smart environments through speculative prototyping and participatory design with building users. By developing interventions that scaffold lived experience, we can shift smart building design from provision and control toward enabling ongoing, meaningful participation in everyday smart buildings.
